\section{Related Work}
\label{sec:related-work}

\subsection{ZK-Friendly Hash Functions}

The dominant hash functions in ZK systems operate natively over prime
fields, avoiding the bitwise operations that make SHA-256 and Keccak
prohibitively expensive in arithmetic circuits.

\textbf{Poseidon}~\cite{grassi2021poseidon} introduced an
algebraic sponge construction optimized for prime-field arithmetic,
achieving 8$\times$ fewer constraints than Pedersen hashes in Groth16
circuits. The design uses partial rounds (applying the S-box to only
one state element) to reduce constraint counts while maintaining
security against algebraic attacks over the BN254 and BLS12-381 fields.

\textbf{Poseidon2}~\cite{grassi2023poseidon2} is a refined variant
that replaces the MDS matrix with a more efficient diffusion layer,
reducing the number of multiplications per round. Benchmarks on the
Goldilocks field report 2--3$\times$ improvement over original
Poseidon in both native and in-circuit performance. Polygon zkEVM's
zkProver uses Poseidon over Goldilocks with 12 state elements and
achieves hashing rates sufficient for processing 1,000+ transactions
per batch~\cite{polygon2024zkprover}.

\textbf{MiMC}~\cite{albrecht2016mimc} uses a simpler $x^3$ or $x^7$
round function but requires more rounds for equivalent security,
resulting in deeper circuits. It remains used in legacy systems but
has been largely superseded by Poseidon in new deployments.

\textbf{Rescue-Prime}~\cite{aly2020rescue} provides an alternative
algebraic hash with stronger security analysis but higher constraint
counts than Poseidon, limiting its adoption in production systems
where proving time is the primary bottleneck.

\subsection{Sparse Merkle Trees}

A Sparse Merkle Tree (SMT) is a fixed-depth binary Merkle tree where
most leaves are empty (containing a default zero value). The key
property is that inclusion and non-inclusion proofs have identical
structure: a Merkle path of length equal to the tree depth,
regardless of the number of populated leaves~\cite{laurie2012revocation}.

\textbf{Depth selection} is a critical design parameter. Ethereum's
state trie addresses are 160 bits (20-byte addresses), while storage
tries use 256-bit keys. Production zkEVM implementations vary:
Polygon zkEVM uses depth 256 with key hashing~\cite{polygon2024zkprover},
zkSync Era uses depth 256 with a custom sparse
tree~\cite{zksync2024boojum}, and Scroll uses a modified MPT
with Poseidon hashing at variable depth~\cite{scroll2024architecture}.

\textbf{Compact variants} such as the Jellyfish Merkle
Tree~\cite{aptos2024jellyfish} (used in Aptos/Diem) employ
internal extension nodes to skip long runs of single-child
paths, reducing both storage and proof size. These optimizations
become critical at depths beyond 64 where the full sparse
structure incurs $O(\text{depth})$ hash computations per update.

\subsection{MPT vs. SMT for EVM State}

The Merkle Patricia Trie (MPT) is Ethereum's native state
structure~\cite{wood2024ethereum}, combining a radix-16 trie with
branch, extension, and leaf nodes. Its advantages include:
dynamic depth (paths compress naturally), Keccak compatibility with
existing EVM infrastructure, and mature tooling (go-ethereum's
\texttt{trie} package).

However, MPT has fundamental disadvantages for ZK proving:
(1)~variable-depth proofs require conditional circuit logic,
increasing constraint count; (2)~Keccak-256 costs approximately
150,000 constraints per hash in R1CS~\cite{grassi2021poseidon}
vs.\ approximately 300 for Poseidon; and (3)~the 16-ary branching
factor means each level requires 16 child hashes, expanding
proof width.

SMT provides fixed-depth proofs (simplifying circuit design),
native compatibility with ZK-friendly hash functions, and
uniform proof structure for both inclusion and exclusion.
The trade-off is higher worst-case proof size at large depths
and the need for compact optimizations when the key space
is sparse relative to $2^{\text{depth}}$.

Production systems have converged on SMT with ZK-friendly hashing.
Polygon zkEVM reports 2--3~ms per Merkle proof verification in
their zkProver~\cite{polygon2024zkprover}. Scroll's modified
zktrie achieves similar performance with a Poseidon-based binary
trie~\cite{scroll2024architecture}.

\subsection{Go Implementations}

The Go ecosystem for ZK-friendly cryptographic primitives has matured
significantly:

\textbf{gnark-crypto}~\cite{gnark2024crypto} (ConsenSys) provides
assembly-optimized BN254 field arithmetic with Poseidon2 support.
The \texttt{fr.Element} type uses Montgomery form with platform-specific
assembly (AMD64, ARM64), avoiding the overhead of Go's
\texttt{math/big.Int}. Benchmarks report Poseidon2 hashing at
3--5~$\mu$s per 2-to-1 compression on modern hardware.

\textbf{vocdoni/arbo}~\cite{vocdoni2024arbo} implements a compact
SMT in Go with Poseidon (original, not Poseidon2) using the
iden3/go-iden3-crypto library. It supports persistent storage via
BadgerDB/Pebble and provides Circom-compatible hashing, making it
directly usable with snarkjs for proof generation.

\textbf{iden3/go-iden3-crypto}~\cite{iden3gocrypto} provides
Go implementations of Poseidon, Baby JubJub, and other
ZK primitives compatible with circomlib's JavaScript
implementations.

\subsection{Production zkEVM State Databases}

Table~\ref{tab:production-systems} summarizes state database
designs in production zkEVM systems.

\begin{table}[h]
\centering
\caption{State database designs in production zkEVM systems}
\label{tab:production-systems}
\begin{tabular}{lllll}
\toprule
\textbf{System} & \textbf{Tree} & \textbf{Hash} & \textbf{Depth} & \textbf{Language} \\
\midrule
Polygon zkEVM & SMT & Poseidon & 256 & C++/Go \\
zkSync Era & SMT & Blake2s & 256 & Rust \\
Scroll & mod.\ MPT & Poseidon & Variable & Go \\
Aztec & Indexed MT & Pedersen & 32 & C++/TS \\
Linea & SMT & MiMC & 40 & Go/Java \\
\bottomrule
\end{tabular}
\end{table}

The convergence toward SMT with algebraic hashing validates the
architectural direction of this work. The primary open question
is the quantitative performance of Go-native implementations at
enterprise scale, which this paper addresses.
